\documentclass[conference]{IEEEtran}
\usepackage[utf8]{inputenc}
\usepackage[T1]{fontenc}
\usepackage{cite}
\usepackage{amsmath,amssymb}
\usepackage{graphicx}
\usepackage{booktabs}
\usepackage{url}

\title{When Validators Don't Trigger: A Paired Evaluation of Generate, Validate, and Repair for LLM‑Generated Network Configurations}

\author{
\IEEEauthorblockN{Autonomous AI Researcher}
\IEEEauthorblockA{CNSM 2026 Agentic AI Researcher Track}
}

\begin{document}

\maketitle

\begin{abstract}
We preregistered and executed a paired confirmatory experiment (N = 600 intent\ensuremath{\rightarrow}configuration task pairs) to test whether a deterministic validator‑guided repair pipeline (generate\ensuremath{\rightarrow}validate\ensuremath{\rightarrow}repair) reduces high‑risk semantic misconfigurations relative to single‑shot LLM generation. The run used shared\_initial\_candidate semantics: exactly one sampled initial candidate per pair was produced and reused by both arms; the guarded arm could differ from baseline only if the deterministic validator rejected that candidate and an automated repair executed. Using gpt‑5‑mini as generator and deterministic\_netops\_validator\_v5 as validator, every shared initial candidate passed the validator in both baseline and guarded conditions (baseline = 600/600; guarded = 600/600). The paired absolute difference = 0.0 (paired bootstrap 95\% CI = [0.0, 0.0]; exact McNemar two‑sided p = 1.0; bootstrap resamples = 10,000, seed = 1729). No automated repairs were invoked. Per the preregistration, a seeded 100‑case validator sensitivity audit was planned but was not executed for this confirmatory run; that omission limits claims about validator construct validity and sensitivity. We report the preregistered design, exact paired results from archived artifacts, clarify the scientific meaning of shared\_initial\_candidate semantics (it measures validator/repair trigger efficacy, not independent generator differences), and give concrete design recommendations for future generate\ensuremath{\rightarrow}validate\ensuremath{\rightarrow}repair evaluations.
\end{abstract}

\section{Introduction}
Translating high‑level operational intents into vendor‑specific device configurations is an active NetOps automation problem with immediate operational risk if translations are semantically incorrect. Prior prototype systems and surveys recommend pairing large language model (LLM) generators with deterministic validators and tool‑augmented repair loops to mitigate hallucinations and semantic errors [1]--[4]. However, how to evaluate such generate\ensuremath{\rightarrow}validate\ensuremath{\rightarrow}repair workflows in a reproducible, hypothesis‑driven way remains unsettled: dataset scope, validator definitions, and generation semantics change what an experiment actually measures. We preregistered and executed a narrowly scoped paired confirmatory experiment to answer a precise estimand under an explicit generation contract. Important methodological point: under shared\_initial\_candidate semantics (one sampled initial generation per task reused by both arms), the paired contrast can only detect validator/repair effects acting on the same sampled candidate; it does not measure independent per‑condition generator variability. Our contributions are: (i) a preregistered paired experiment executed on 600 intent\ensuremath{\rightarrow}configuration pairs with archived artifacts; (ii) an artifact‑grounded report of the null (ceiling) outcome that arises under shared\_initial\_candidate semantics; and (iii) concrete recommendations (sampling semantics, mandatory validator audits, adversarial stress) for future evaluations of generate\ensuremath{\rightarrow}validate\ensuremath{\rightarrow}repair pipelines.

\section{Related Work}
Work on intent\ensuremath{\rightarrow}configuration translation and LLM‑assisted NetOps spans prototype systems and production reports. Intent‑translation prototypes show feasibility of mapping natural language intents to vendor configurations [1]; a production multi‑agent framework integrating validators and tool use (Confucius) documents large‑scale operational integration and emphasizes validator machinery in practice [2]. Recent surveys argue that agentic NetOps and AIOps reliability depends more on surrounding tooling, validators, and governance than on the base LLM alone [3]. Separate literature on LLM self‑correction and tool augmentation shows that tool access or external calculators can substantially reduce specific error classes in other domains (e.g., clinical calculations), though gains depend on tool quality and integration [4], [7]. Surveys of hallucination and of constrained/controlled decoding document the failure modes we aim to exercise and outline techniques (constrained decoding, neurosymbolic repair) that can enforce syntactic or richer constraints at generation time [5], [6]. These works motivated our preregistered paired confirmatory design with explicit generation semantics and a deterministic validator; we do not generalize empirical improvement claims beyond the domains covered by the cited studies.

\section{Methodology}
Preregistration and estimand: The study was preregistered as autocand-2026-01-prereg-v1. The primary estimand is the paired absolute difference in high‑risk semantic failure proportions (guarded - baseline) across a prereclared set of N = 600 pairs. The analysis plan used an exact McNemar two‑sided test and paired bootstrap 95\% CIs (10,000 resamples, seed = 1729), executed by the archived paired\_binary\_analysis\_v1 executor. Sampling and generator contract: tasks (task\_count = 600) were deterministically selected via the adapter's challenging\_workflow\_stress\_v1 rule; planned episodes = 1,200. The adapter contract specified generation\_semantics = shared\_initial\_candidate: exactly one sampled initial generation per task was produced and reused unchanged by both baseline and guarded conditions. Per contract: initial\_generation\_calls\_per\_task = 1 and maximum\_repair\_calls\_per\_task = 1; the guarded arm could invoke deterministic\_netops\_validator\_v5 and request at most one automated repair only if that validator rejected the shared initial candidate. Model and runtime metadata: the declared generator was gpt‑5‑mini and the archived model configuration records temperature = null/unset; null/unset is not evidence of deterministic hosted‑model sampling and does not imply temperature = 0. Validator audit (preregistered but unexecuted): the preregistration included a seeded 100‑case validator sensitivity/specificity audit to assess validator construct validity; that audit was not executed for this confirmatory run and therefore no independent sensitivity estimate for deterministic\_netops\_validator\_v5 is reported here. Inclusion/exclusion and contamination controls followed the preregistered rules (public‑source requirement, provenance checks, algorithmic redaction of unsafe exemplars). Analysis executor: paired\_binary\_analysis\_v1 produced the confirmatory test and bootstrap CIs using the archived result set (bootstrap resamples = 10,000, seed = 1729).

\section{Results}
Execution accounting and primary outcomes (archived): the preregistered run completed to plan. planned\_episode\_count = 1,200; completed\_episode\_count = 1,200; complete\_pair\_count = 600; model\_calls\_used = 600 (one shared initial generation per pair under shared\_initial\_candidate semantics); failed episodes = 0. Confirmatory outcomes from archived analysis artifacts: baseline\_success\_count = 600/600 (100.0\%); guarded\_success\_count = 600/600 (100.0\%). Paired contingency counts: n11 = 600, n10 = 0, n01 = 0, n00 = 0. The paired estimate (guarded - baseline) = 0.0. Paired bootstrap 95\% CI (percentile, 10,000 resamples, seed = 1729) = [0.0, 0.0]. Exact McNemar two‑sided p = 1.0. Repair attempts invoked: none --- deterministic\_netops\_validator\_v5 accepted every shared initial candidate so no automated repairs were applied. Planned validator audit not executed: as prespecified, the seeded 100‑case validator sensitivity audit was not run; therefore no confirmatory inference about validator sensitivity or false permissiveness is available from this execution and the possibility of an overpermissive validator producing a ceiling effect remains an unresolved empirical limitation.

\section{Discussion}
Scientific meaning of the observed null outcome: under shared\_initial\_candidate semantics the guarded pipeline can affect a paired outcome only if the deterministic validator rejects the same sampled initial candidate and a repair produces a validator‑passing configuration. Because the validator accepted all initial candidates, the guarded pipeline had no opportunity to show benefit; the observed ceiling outcome documents lack of validator rejections on this task bank rather than generator superiority. Design implications: (1) align generation semantics to the estimand --- use independent per‑condition sampling to measure generator‑level interventions and shared\_initial\_candidate semantics when measuring validator/repair trigger efficacy; (2) preregister and execute validator sensitivity audits (seeded adversarial perturbations or held‑out failure cases) before confirmatory tests to detect overpermissive validators that create ceiling effects; (3) include adversarially stressed or harder task partitions to increase the likelihood of validator triggers and repair opportunities. Power note: the preregistration power calculation assumed a nonzero baseline failure rate (\textasciitilde{}10\%); the observed baseline failure rate was 0\%, so the preregistered minimal detectable effect (5 percentage points) and intended power interpretation do not apply to this execution. Scope and reproducibility: tasks were synthetic items from deterministic\_netops\_task\_generator\_v5, only one generator (gpt‑5‑mini) and one validator implementation were used, and the planned 100‑case validator audit was not executed; these constraints limit external generality and prevent concluding about validator construct validity beyond the archived pass/fail outcomes.

\section{Limitations}
\begin{itemize}
\item Shared\_initial\_candidate semantics: baseline and guarded used the identical sampled initial generation per pair; thus the paired estimand measures validator/repair trigger effectiveness, not independent generator differences.
\item Validator construct validity: deterministic\_netops\_validator\_v5 accepted all initial candidates in this run; without the preregistered seeded 100‑case validator audit (not executed) an independent sensitivity/specificity estimate is unavailable, creating the risk of an overpermissive validator producing a ceiling effect.
\item Task external validity: tasks were synthetic items generated by deterministic\_netops\_task\_generator\_v5 under the preregistered contamination and redaction rules; results may not generalize to diverse production device configurations or vendor idiosyncrasies.
\item Model and adapter scope: only the declared model (gpt-5-mini) and one adapter/validator instance were used; no claim of model‑class generality is made.
\item Sampling record: the archived model configuration records temperature = null/unset; null/unset is explicitly not evidence of deterministic sampling. Independent per‑condition sampling would be required to evaluate generator‑level effects.
\item Power assumptions violated: the preregistration's power calculation assumed a nonzero baseline failure rate (\textasciitilde{}10\%). The observed baseline failure rate was 0\%, so the preregistered minimal detectable effect (5 percentage points) and intended power interpretation do not apply to this execution.
\item Planned validator audit not executed: the preregistered seeded 100‑case validator audit was not executed as part of the archived confirmatory run; validator sensitivity therefore remains an unresolved limitation and a concrete requirement for follow‑up work.
\end{itemize}

\section{Conclusion}
We executed a preregistered paired experiment (N = 600 pairs) under shared\_initial\_candidate semantics comparing single‑shot LLM generation and a deterministic validator‑guided repair pipeline. Both arms passed the deterministic validator in all pairs (100\% baseline and guarded), yielding an observed paired difference of 0.0 (95\% CI [0.0, 0.0], p = 1.0). No automated repairs were invoked. The result highlights that (a) generation semantics determine what a paired estimand measures (validator trigger effectiveness vs. generator differences), (b) validators must be audited to avoid ceiling effects that mask downstream protections, and (c) future experiments should use independent per‑condition sampling or adversarial/harder task banks when the scientific goal is to assess generator‑level interventions. All preregistration, execution, validator traces, and analysis artifacts referenced in the preregistration are archived per the preregistered plan.

\section*{Disclosure Statement}
Models, tools, and preregistration: generative calls used the declared model gpt-5-mini via the hosted\_netops\_gvr\_v1 adapter; the deterministic validator was deterministic\_netops\_validator\_v5; tasks were produced by deterministic\_netops\_task\_generator\_v5. The study was preregistered as autocand-2026-01-prereg-v1 and executed under that preregistration. Autonomy boundary: a human Research Director selected the broad topic family and approved the initial master prompt before the autonomy lock; after the lock the AI system autonomously selected the final research question, executed the preregistered experiment, performed analyses, and authored and revised the manuscript. No manual human editing of technical content, numerical results, or manuscript conclusions occurred after the autonomy lock. Master prompt reference (immutable): provenance/master\_prompt.txt, SHA‑256 = 1872df1e\allowbreak{}1805d2d9\allowbreak{}6940456c\allowbreak{}a016bd66\allowbreak{}5d1d5196\allowbreak{}add77f5a\allowbreak{}cdf1582b\allowbreak{}b39b15ba. Generation semantics and sampling: generation\_semantics = shared\_initial\_candidate (one sampled initial generation per task was reused unchanged by baseline and guarded); the archived model configuration records temperature = null/unset for executed calls --- null/unset is not evidence of deterministic sampling and does not imply temperature = 0. The preregistered seeded 100‑case validator sensitivity audit was not executed for this confirmatory run; that omission is reported transparently in Methods and Results and limits claims about validator construct validity.

\begin{thebibliography}{99}
\bibitem{ref1} Nguyen Tu, Sukhyun Nam, James Won‐Ki Hong, "Intent‐Based Network Configuration Using Large Language Models", 2024, 10.1002/nem.2313
\bibitem{ref2} Zhaodong Wang, Samuel Lin, Guanqing Yan, Soudeh Ghorbani, Minlan Yu, Jiawei Zhou, Nathan Hu, Lopa Baruah, Sam Peters, Srikanth Kamath, Jian Yang, Ying Zhang, "Intent-Driven Network Management with Multi-Agent LLMs: The Confucius Framework", 2025, 10.1145/3718958.3750537
\bibitem{ref3} Muhammad Bilal, Jon Crowcroft, Ruizhi Wang, Xiaolong Xu, Schahram Dustdar, "Large Language Models for Agentic NetOps and AIOps: Architectures, Evaluation, and Safety", 2026, 10.48550/arxiv.2605.12729
\bibitem{ref4} Liangming Pan, Michael Saxon, Wenda Xu, Deepak Nathani, Xinyi Wang, William Yang Wang, "Automatically Correcting Large Language Models: Surveying the landscape of diverse self-correction strategies", 2023, 10.48550/arxiv.2308.03188
\bibitem{ref5} Yue Zhang, Yafu Li, Leyang Cui, Cai Deng, Lemao Liu, Tingchen Fu, Xinting Huang, Enbo Zhao, Yanwen Zhang, Yulong Chen, Longyue Wang, Ahn Tuan Luu, Wei Bi, Freda Shi, Shuming Shi, "🧜Siren's Song in the AI Ocean: A Survey on Hallucination in Large Language Models", 2025, 10.1162/coli.a.16
\bibitem{ref6} Patrick Pynadath, Ruqi Zhang, "Controlled LLM Decoding via Discrete Auto-regressive Biasing", 2025, 10.48550/arxiv.2502.03685
\bibitem{ref7} Alex J. Goodell, Simon N. Chu, Dara Rouholiman, Larry F. Chu, "Large language model agents can use tools to perform clinical calculations", 2025, 10.1038/s41746-025-01475-8
\bibitem{ref8} Wayne Xin Zhao, Kun Zhou, Junyi Li, Tianyi Tang, Zican Dong, Yupeng Hou, Beichen Zhang, Yingqian Min, Junjie Zhang, Peiyu Liu, Xiaolei Wang, Yifan Du, Yushuo Chen, Yushuo Chen, Zhipeng Chen, Jinhao Jiang, Ruiyang Ren, Yifan Li, Xinyu Tang, Peiyu Liu, Yiwen Hu, Jian‐Yun Nie, Ji-Rong Wen, "A Survey of Large Language Models", 2026, 10.1007/s11704-026-60308-3
\end{thebibliography}

\end{document}
